\section{Conclusion}

Le projet AGRI-Scout aura été, avant tout, une leçon sur la distance qui sépare la conception théorique de la réalité du déploiement matériel. Les obstacles rencontrés (instabilité électrique, synchronisation Python/Arduino, contraintes d'impression 3D, limitations de calcul pour la vision) ont conduit l'équipe à repenser ses choix de façon itérative, aboutissant souvent à des solutions plus simples mais plus robustes que celles envisagées initialement. Le prototype final remplit ses objectifs fondamentaux : navigation autonome par suivi de ligne, détection d'obstacles par LiDAR, mécanisme de sondage du sol conçu et fabriqué en 3D, et collecte automatisée de données multiparamétriques via la sonde industrielle RS485.

Ce projet a permis à l'équipe d'acquérir une expérience pratique dans des domaines rarement enseignés en cours magistral : gestion des conflits électroniques en système embarqué, conception mécanique itérative, débogage de systèmes distribués en temps réel. La valeur de ce travail ne réside pas uniquement dans le robot livré, mais aussi dans la documentation des impasses techniques. Savoir ce qui n'a pas fonctionné est essentiel. Les solutions abandonnées et les pivots stratégiques constituent, pour les futurs développeurs, un guide pour éviter de répéter les mêmes erreurs et accélérer la prise en main du système.

L'AGRI-Scout est un prototype qui peut encore progresser. La réintégration de la détection YOLO pour la reconnaissance des végétaux et des pestes, la navigation par waypoints GPS RTK, et l'utilisation d'une plateforme de calcul plus puissante sont des pistes concrètes pour les prochaines équipes qui reprendront ce projet.
